La construcción de una interfaz para gestionar y corregir exámenes exige
coordinar formularios, navegación, documentos y operaciones remotas. Conviene
distinguir la biblioteca que representa las vistas, las herramientas que generan
el artefacto distribuible y los mecanismos que permiten verificar su comportamiento.
Esta sección presenta esos niveles y los relaciona con decisiones observables
en SCDE. No constituye una comparación experimental de rendimiento entre
tecnologías ni presupone que una biblioteca garantice la calidad de la aplicación.

\subsection{Componentes y representación del estado}

React propone descomponer una interfaz en componentes, describir sus estados
visuales y conectar el flujo de datos entre ellos. Su documentación recomienda
conservar una representación mínima del estado: los valores calculables a partir
de otros datos no necesitan almacenarse de nuevo. Cuando varios componentes
comparten información mutable, su propiedad puede situarse en un antecesor común,
que transmite datos y funciones de actualización mediante propiedades.
Este enfoque permite razonar de forma explícita sobre qué información determina
una vista y qué interacción puede modificarla~\cite{codex_react_thinking}.

En una aplicación administrativa, esta distinción ayuda a separar los registros
recibidos del servicio, los valores aún no enviados de un formulario y los
criterios locales de presentación. Por ejemplo, una búsqueda limitada a la página
actual no representa una consulta global al servidor. Esa diferencia pertenece al
significado de los datos y debe conservarse tanto en el estado como en el texto
de la interfaz. La composición de componentes facilita expresar estas fronteras,
pero su definición sigue siendo una responsabilidad del diseño de la aplicación.

La base inspeccionada de SCDE declara React~19 y utiliza módulos JavaScript con
JSX. La organización por funcionalidades añade una convención propia: cada módulo
expone rutas, navegación, traducciones y simulaciones mediante un registro común.
Esta estructura es una decisión del repositorio, no una exigencia de React.
Del mismo modo, el uso de JavaScript es un hecho verificable de la implementación;
no permite afirmar una ventaja medida sobre TypeScript. La ausencia de tipos
estáticos propios exige prestar especial atención a las formas de los datos en
los límites de comunicación~\cite{codex_scde_frontend_architecture}.

\subsection{Navegación y separación de responsabilidades}

React Router diferencia varios modos de utilización. El modo declarativo ofrece
la asociación entre direcciones y componentes, navegación y estados de enlaces.
El modo de datos añade mecanismos para cargas y acciones, mientras que el modo
de \textit{framework} incorpora convenciones adicionales de construcción y
representación. Esta distinción evita atribuir a cualquier uso de la biblioteca
las capacidades de todos sus modos~\cite{codex_router_modes}.

En SCDE, las funcionalidades declaran sus rutas y el contenedor común aplica
guardas de acceso. Las peticiones se encapsulan en funciones de API y los estados
de carga y acción se gestionan mediante utilidades compartidas. Por tanto, la
arquitectura observada separa navegación y acceso a datos, sin que la dependencia
de React Router pruebe por sí sola el uso de sus cargadores o de representación
en servidor~\cite{codex_scde_frontend_architecture}.
Las guardas mejoran la coherencia de la navegación, pero la autorización de una
operación debe comprobarse en el servicio que la ejecuta. Ocultar un enlace no
constituye una garantía de autorización.

\subsection{Construcción y distribución de la aplicación}

La documentación de Vite~7 describe una construcción de producción a partir de
una entrada HTML, destinada a generar recursos servibles desde alojamiento
estático. También permite configurar la ruta base cuando la aplicación se sirve
desde un subdirectorio. De esta manera, el proceso de construcción transforma
fuentes y dependencias en un artefacto que puede distribuirse por separado del
entorno de desarrollo~\cite{codex_vite7_build}.

El repositorio declara Vite~7 y configura su integración con React. La aplicación
puede compilarse para usar el cliente simulado o el servicio real mediante su
configuración. La generación correcta de recursos acredita que el proceso de
construcción termina; no acredita la disponibilidad del backend, la validez de una
sesión ni el funcionamiento de todas las operaciones. En una aplicación con
navegación de cliente también debe comprobarse el acceso directo a direcciones
internas desde el alojamiento elegido. Esas comprobaciones de distribución son
distintas de las pruebas de componentes~\cite{codex_scde_frontend_architecture}.

\subsection{Verificación de interacciones y límites del entorno simulado}

Testing Library orienta sus utilidades hacia el uso observable de los componentes:
favorece trabajar con nodos del DOM y con interacciones cercanas a las que realiza
una persona, en lugar de depender de instancias internas. Este criterio ayuda a
formular comprobaciones sobre acciones y resultados visibles y a reducir el
acoplamiento de una prueba con detalles internos que pueden cambiar sin alterar
el comportamiento~\cite{codex_testing_principles}.

SCDE combina Vitest, Testing Library y un entorno jsdom configurado en Vite. Las
pruebas utilizan el modo simulado del cliente y se complementan con ESLint,
compilación y una auditoría de estilos. Son comprobaciones con objetivos
diferentes: detectar incumplimientos estáticos, comprobar interacciones y producir
el artefacto. Su combinación aporta evidencia sobre los casos ejecutados, sin
convertir la simulación del DOM en una ejecución completa de navegador ni los
datos simulados en una validación del backend real.
La corrección de contratos, la accesibilidad y la experiencia de uso necesitan
evidencias adicionales adecuadas a cada propiedad~\cite{codex_scde_frontend_architecture}.

La selección tecnológica permite así distribuir responsabilidades entre vistas,
navegación, construcción y verificación. Las funciones de internacionalización,
presentación de PDF y primitivas visuales completan la base del proyecto, pero
su mera inclusión en las dependencias no demuestra traducción completa,
accesibilidad documental ni consistencia de todas las pantallas. La valoración
de la solución debe apoyarse en la implementación y en los resultados registrados,
además de en las capacidades documentadas de las herramientas.
